% =====================================================================
% A/B results: live-measured comparison of Environment A (centralized
% baseline) vs Environment B (compartmentalized DNAT). English, LADC/acmart.
% Self-contained: drop-in tables + a Results subsection that complements
% section_evaluation.tex. Requires booktabs + tabularx (already in preamble).
% Data: TCC/evaluation_data/results_ab_campaign.md (measured 2026-06-11).
% =====================================================================

% ---------------------------------------------------------------------
% Subsection prose: paste at the end of the Evaluation section (or just
% before "Threats to Validity"), after the A column has been promoted
% from "derived from configuration" to "measured live".
% ---------------------------------------------------------------------
\subsection{A/B Comparison: Live Measurement of Both Environments}
\label{sec:ab-results}

To remove any inference from the comparison, we ran the full suite against both deployments
live and end to end: the compartmentalized system (Environment~B) and the centralized
baseline (Environment~A, a single container colocating the same logic). Because both share
the same Firecracker microVMs and the same access gate, the per-microVM guarantees are a
constant---T1, T1b, T3 and T5 reproduce in both, and so does the operational cost---so the
architectural variable surfaces only in the blast radius. That is where A and B diverge
(Table~\ref{tab:blast-radius}, populated from live inspection rather than configuration):
in B, the executor and the builder carry \emph{none} of the control-plane secrets, and the
sensitive stores (IPFS repository, wheel cache, executions store) are absent from their
namespace, whereas in A those same secrets and stores sit in the single namespace that also
runs the execution logic. The supply-chain test (T2) inherits the split: a compromised
dependency's code executes, but in B inside a plane that holds neither secrets nor a route
to them.

\textbf{What the network residual still grants.} Reachability is not exploitability, so we
state what a compromised executor can \emph{do} with the two ports it still reaches in B.
The prototype does not authenticate its internal endpoints, so the control API can be driven
as a confused deputy: an attacker on the executor can enumerate registered assets, read the
stored output of previous executions, and ask the control plane to act with its wallet
through the fixed set of operations the API exposes. This is a real residual, and it is why
internal mutual TLS heads our future work. What it does \emph{not} give is the platform: the
encryption and wallet keys are never returned by any endpoint and remain in CVM1's
environment, so the attacker can neither decrypt artifacts pulled from the store nor sign
transactions of its own choosing, and the on-chain gate still refuses an execution for which
no access was purchased (T5). In A the same attacker needs no network hop at all---both keys
sit in its own environment and the stores on its filesystem---so it can decrypt every
archived artifact and sign arbitrary transactions directly. The difference is therefore not
the presence of a residual but its ceiling: a bounded, mediated set of API calls in B
against unmediated custody of the platform's keys in A---\emph{containment, not absolute
security}.

% ---------------------------------------------------------------------
% Table: blast radius B vs A — English version, live-measured wording.
% A/B cost figures live in tab:comparison (section_evaluation.tex).
% ---------------------------------------------------------------------
\begin{table}[htbp]
\centering
\caption{Reach of a compromised execution plane over control-plane assets: the
compartmentalized design (B) versus the centralized baseline (A), from live inspection.
``Isolated'' = absent from the plane's namespace (the environment variable is not set, the
store is not mounted); ``Colocated'' = same namespace as the execution logic.}
\label{tab:blast-radius}
\footnotesize
\begin{tabularx}{\textwidth}{@{}>{\raggedright\arraybackslash}X>{\raggedright\arraybackslash}p{2.4cm}>{\raggedright\arraybackslash}p{2.4cm}@{}}
\toprule
Control-plane asset reachable from the execution plane & Env.~B (proposed) & Env.~A (baseline) \\
\midrule
Asset-encryption key (env) & Isolated & Colocated \\
Wallet private key (env) & Isolated & Colocated \\
Service URLs: RPC, IPFS, builder, executor (env) & Isolated & Colocated \\
IPFS repository, wheel cache, executions store (filesystem) & Isolated & Colocated \\
IPFS API (port 5001, network) & Blocked & Reachable \\
Control API (3001) and blockchain RPC (8545), network & Reachable$^{\dagger}$ & Reachable \\
\bottomrule
\end{tabularx}
\par\smallskip
\parbox{\textwidth}{\footnotesize $^{\dagger}$Residual risk in the single-host deployment:
the three containers share one Docker network and the internal endpoints are
unauthenticated, so a compromised executor can invoke the control API as a confused deputy
(enumerate assets, read stored execution outputs) without ever holding its keys---see the
discussion above. The defining assets of the blast radius---secrets and persistent
stores---are nonetheless isolated in B and colocated in A.}
\end{table}

% ---------------------------------------------------------------------
% Threats to Validity: movida de section_evaluation.tex (texto intacto)
% para fechar a secao de Evaluation depois da subsecao A/B.
% ---------------------------------------------------------------------
\subsection{Threats to Validity}

The measurements were taken on a single host; absolute latencies will vary with hardware (a
colder first campaign saw a 31.4~s execution mean and 103--131~s builds, with A and B again
indistinguishable), and each cost figure is a single run, so we read them as orders of
magnitude rather than precise separations. The containment and exfiltration results, by
contrast, are structural---absence of a network interface, absence of secrets---and do not
depend on timing. Our workload is also deliberately small: a 60-row dataset isolates the
fixed microVM cost, but says nothing about how that cost amortizes over a long-running job,
nor about where the per-run caps (1~vCPU, 512~MiB, fixed disk sizes) begin to bind. The
network-reach residual bounds the blast-radius claim, with the consequences and limits set
out above, and confidentiality against the infrastructure operator (the second half of R1)
is assumed rather than shown, for the substrate reason given next. Finally, the DoS analysis
covers only per-run state hygiene, and the SGX column of Table~\ref{tab:comparison} is
qualitative rather than re-measured.

A more fundamental caveat concerns the execution substrate. The prototype runs each plane as
a container with Firecracker microVMs rather than inside a hardware-backed confidential VM,
so our measurements do not come from a genuinely confidential environment and we make no
claim about confidential-computing overheads. This does not weaken the architectural
argument: the properties we evaluate---privilege separation across planes, a networkless
execution microVM, ephemeral teardown, and on-chain access control---belong to the
architecture, not to the memory-encryption substrate. Running the same planes inside
confidential VMs (e.g., AMD SEV-SNP~\cite{amd2020sevsnp}) would add attestation and
memory-encryption costs---so the cost figures above would not carry over unchanged---but
would preserve the containment and blast-radius results, which is why we treat confidential
execution as a substrate substitution rather than an architectural change.
